\subsection{TENSOR PRODUCT OF TWO MODULES OVER A COMMUTATIVE RING}

Let C be a commutative ring; for every C-module E, the module structure on E is compatible with itself (§ 1, no. 14). If E and F are two C-modules, the considerations of no. 4 then allow us to define two C-module structures on the tensor product E ⊗c F, respectively such that γ(x ⊗ y) = (γx) ⊗ y and such that γ(x ⊗ y) = x ⊗ (γy); but as, by Definition 1 of no. 1, in this case (γx) ⊗ y = x ⊗ (γy), these two structures are the same. Henceforth when we speak of E ⊗c F as a C-module, we shall mean with the structure thus defined, unless otherwise stated. The canonical isomorphism.

\[
\sigma : \mathbf {F} \otimes_ {\mathbf {c}} \mathbf {E} \rightarrow \mathbf {E} \otimes_ {\mathbf {c}} \mathbf {F}
\]

(no. 1, Corollary 2 to Proposition 1) is then a C-module isomorphism.

It follows from this definition that, if \((a_{\lambda})_{\lambda\in\mathbb{L}}\) (resp. \((b_{\mu})_{\mu\in\mathbb{M}}\) ) is a generating system of the C-module E (resp. F), \((a_{\lambda}\otimes b_{\mu})\) is a generating system of the C-module \(E\otimes_{C}F\) ; in particular, if E and F are finitely generated C-modules, so is \(E\otimes_{C}F\) .

For every C-module G, the Z-bilinear mappings \(f\) of \(\mathbf{E} \times \mathbf{F}\) into \(\mathbf{G}\) for which

\[
f (\gamma x, y) = f (x, \gamma y) = \gamma f (x, y) \quad \text { for } x \in \mathbf {E}, y \in \mathbf {F}, \gamma \in \mathbf {C}\tag{9}
\]

are then called C-bilinear and form a C-module denoted by \(\mathcal{L}_2(\mathrm{E},\mathrm{F};\mathrm{G})\) ; Proposition 3 (no. 4) defines a canonical C-module isomorphism (cf.~§1, no. 14, Remark 1).

\[
\mathcal {L} _ {2} (\mathrm{E}, \mathrm{F}; \mathrm{G}) \rightarrow \operatorname{Hom} _ {\mathbf {c}} (\mathrm{E} \otimes_ {\mathbf {c}} \mathrm{F}, \mathrm{G}).\tag{10}
\]

Let \(E'\) , \(F'\) be two C-modules and \(u: E \to E'\) , \(v: F \to F'\) two C-linear mappings; then (no. 4) \(u \otimes v\) is a C-linear mapping of \(E \otimes_{\mathbb{C}} F\) into \(E' \otimes_{\mathbb{C}} F'\) . Further, it is immediate that \((u, v) \mapsto u \otimes v\) is a C-bilinear mapping of \(\operatorname{Hom}_{\mathbb{C}}(\mathbf{E}, \mathbf{E}') \times \operatorname{Hom}_{\mathbb{C}}(\mathbf{F}, \mathbf{F}')\) into \(\operatorname{Hom}_{\mathbb{C}}(\mathbf{E} \otimes_{\mathbb{C}} \mathbf{F}, \mathbf{E}' \otimes_{\mathbb{C}} \mathbf{F}')\) ; hence there canonically corresponds to it a C-linear mapping, called canonical:

\[
\operatorname{Hom} _ {\mathbf {c}} (\mathbf {E}, \mathbf {E} ^ {\prime}) \otimes_ {\mathbf {c}} \operatorname{Hom} _ {\mathbf {c}} (\mathbf {F}, \mathbf {F} ^ {\prime}) \to \operatorname{Hom} _ {\mathbf {c}} (\mathbf {E} \otimes_ {\mathbf {c}} \mathbf {F}, \mathbf {E} ^ {\prime} \otimes_ {\mathbf {c}} \mathbf {F} ^ {\prime})\tag{11}
\]

which associates with every element \(u \otimes v\) of the tensor product

\[
\operatorname{Hom} _ {\mathbf {c}} (\mathbf {E}, \mathbf {E} ^ {\prime}) \otimes_ {\mathbf {c}} \operatorname{Hom} _ {\mathbf {c}} (\mathbf {F}, \mathbf {F} ^ {\prime})
\]

the linear mapping \(u \otimes v\) . Note that the canonical mapping (11) is not necessarily injective nor surjective (§ 4, Exercise 2).

Remarks. (1) Let A, B be two commutative rings, \(\rho:B\to A\) a ring homomorphism and E and F two A-modules; then the canonical mapping (6) of no. 3 is a B-linear mapping

\[
\rho_ {*} (\mathbf {E}) \otimes \rho_ {*} (\mathbf {F}) \rightarrow \rho_ {*} (\mathbf {E} \otimes_ {\mathbf {A}} \mathbf {F}).\tag{12}
\]

\begin{enumerate}
\def\labelenumi{(\arabic{enumi})}
\setcounter{enumi}{1}
\tightlist
\item
  What has been said in this no. can be generalized to the following case: let E be a right A-module, F a left A-module, C a commutative ring and \(\rho:C\to A\) a homomorphism of C into A such that \(\rho(\mathbf{C})\) is contained in the centre of A (cf.~III, §1, no. 3). We may then consider the C-modules \(\rho_{*}(\mathbf{E})\) and \(\rho_{*}(\mathbf{F})\) and the hypothesis on \(\rho\) implies that these C-module structures are compatible respectively with the A-module structures on E and F (§1, no. 14). The tensor product \(E\otimes_{A}F\) is thus (by virtue of no. 4) given two C-module structures such that \(\gamma(x\otimes y)=(x\rho(\gamma))\otimes y\) and
\end{enumerate}

\[
\gamma (x \otimes y) = x \otimes (\rho (\gamma) y)
\]

respectively for \(\gamma\in\mathbf{C}\) , \(x\in\mathbf{E}\) , \(y\in\mathbf{F}\) and Definition 1 (no. 1) shows also that these two structures are identical. If \(\mathbf{E}'\) (resp. \(\mathbf{F}'\) ) is a right (resp. left) A-module and \(u:\mathbf{E}\to\mathbf{E}'\) , \(v:\mathbf{F}\to\mathbf{F}'\) are two A-linear mappings, then \(u\otimes v:\mathbf{E}\otimes_{\mathbf{A}}\mathbf{F}\to\mathbf{E}'\otimes_{\mathbf{A}}\mathbf{F}'\) is C-linear for the C-module structures just defined; the mapping \((u,v)\mapsto u\otimes v\) :

\[
\operatorname{Hom} _ {\mathbf {A}} (\mathrm{E}, \mathrm{E} ^ {\prime}) \times \operatorname{Hom} _ {\mathbf {A}} (\mathrm{F}, \mathrm{F} ^ {\prime}) \rightarrow \operatorname{Hom} _ {\mathbf {C}} (\mathrm{E} \otimes_ {\mathbf {A}} \mathrm{F}, \mathrm{E} ^ {\prime} \otimes_ {\mathbf {A}} \mathrm{F} ^ {\prime})
\]

is C-bilinear (for the C-module structures on \(\mathrm{Hom}_{\mathbf{A}}(\mathbf{E},\mathbf{E}^{\prime})\) and \(\mathrm{Hom}_{\mathbf{A}}(\mathbf{F},\mathbf{F}^{\prime})\)

defined in § 1, no. 14, Remark 1), whence we also derive a C-linear mapping, called canonical

\[
\operatorname{Hom} _ {\mathbf {A}} (\mathbf {E}, \mathbf {E} ^ {\prime}) \otimes_ {\mathbf {C}} \operatorname{Hom} _ {\mathbf {A}} (\mathbf {F}, \mathbf {F} ^ {\prime}) \rightarrow \operatorname{Hom} _ {\mathbf {C}} (\mathbf {E} \otimes_ {\mathbf {A}} \mathbf {F}, \mathbf {E} ^ {\prime} \otimes_ {\mathbf {A}} \mathbf {F} ^ {\prime}).\tag{13}
\]

\begin{enumerate}
\def\labelenumi{(\arabic{enumi})}
\setcounter{enumi}{2}
\tightlist
\item
  Let \(\mathbf{A}\) be an integral domain and \(\mathbf{K}\) its field of fractions. If \(\mathbf{E}\) and \(\mathbf{F}\) are two vector \(\mathbf{K}\) -spaces, the canonical mapping
\end{enumerate}

\[
\left(\mathbf {E} _ {[ \mathbf {A} ]}\right) \otimes_ {\mathbf {A}} \left(\mathbf {F} _ {[ \mathbf {A} ]}\right)\rightarrow \mathbf {E} \otimes_ {\mathbf {K}} \mathbf {F}
\]

(no. 3 and §1, no. 13) is bijective. It suffices (no. 4) to prove that if \(f\) is an A-bilinear mapping of \(\mathbf{E} \times \mathbf{F}\) into a vector K-space \(\mathbf{G}\) , \(f\) is also K-bilinear. Now, for all \(\alpha \neq 0\) in \(\mathbf{A}\) .

\[
\alpha f (\alpha^ {- 1} x, y) = f (x, y) = \alpha f (x, \alpha^ {- 1} y)
\]

whence

\[
f (\alpha^ {- 1} x, y) = f (x, \alpha^ {- 1} y) = \alpha^ {- 1} f (x, y)
\]

since G is a vector K-space.

